\scanpage{35}
\mspage{33}
% archive sheet number: 34
présent exposé en se réduisant cf. \unsure{6.6|8.6}. à \deleted{\unsure{un cas}} une situation avec un corps résiduel fini. \inserted{Cette démonstration \deleted{\unsure{s'applique pour}} donne la conj. dans les cas a) et c).}
\deleted{Dans} La démonstration géométrique fera l'objet de l'exposé suivant; elle permettra de préciser assez considérablement \deleted{l'énoncé et} la conclusion de unipotence, et \deleted{\unsure{permet de}} s'applique à l'heure actuelle dans les cas a) et b).

7.3. On peut donner une forme précisée de 7.1, en y faisant varier $\ell$ parmi tous les nbs premiers distincts de $p$. Serre et Tate \deleted{se bornent à} demandent que les \deleted{représentations obtenues} groupes $\Gamma$ \deleted{\unsure{opérant dans}} (quotients de $I$) associés à \deleted{\unsure{deux repr. distinctes}} \inserted{$H^i$ pour deux $\ell$ distincts} soient les mêmes, et que les représentations obtenues de $\Gamma$ soient ``isomorphes'' (sur une \unsure{sur-extension} commune $\unsure{\Omega}$ de $\unsure{\mathbb Q_\ell|Q_\ell},\mathbb Q_{\ell'}$); \deleted{\unsure{et}} \inserted{Plus précisément, on peut demander} \deleted{\unsure{la relation}} \inserted{$\operatorname{Tr}g_{H^i(\overline X,\mathbb Q_\ell)}$ soit un entier naturel pour $\forall\,g\in I$, et qu'il soit indépendant de $\ell\ne p$:}
\deleted{\unsure{s'exprime}} \unsure{plus simplement}
\[
 \operatorname{Tr}g_{H^i(\overline X,\mathbb Q_\ell)}
 =\operatorname{Tr}g_{H^i(\overline X,\mathbb Q_{\ell'})}
 \tag*{(7.3.1)}
\]
pour tout $g\in\pi$. [Ceci sera prouvé pour $i=1$ par la suite, en utilisant les modèles de Néron]. \deleted{\unsure{De plus, on est conduit à une question plus considérable}.}

7.4. Lorsque $k$ est fini, \deleted{Serre et Tate} il y a lieu (avec Serre et Tate) de conjecturer que cette \deleted{égalité} \inserted{relation (7.3.1)} est vraie \inserted{(ds $\mathbb Q$)} pour $\forall\,g\in\pi$ dont l'image dans $\operatorname{Gal}(\overline k/k)\simeq\widehat{\mathbb Z}$ tombe dans $\mathbb Z$. [Cela même \deleted{\unsure{pour}} \inserted{sera prouvé} pour $i=1$ dans notre séminaire.]

On veut de plus que si $g\in\pi$ se \unsure{projette} sur \deleted{\unsure{f}} \inserted{$f^{-1}$} (f le générateur
% margin note written slanted
\margin{ou frobenius) les \uuline{valeurs} propres de $g$ opérant sur $H^i(\overline X,\mathbb Q_\ell)$ soient de \unsure{valeur|valeurs} \unsure{absolue complexe|absolues complexes} de la forme $q^{\unsure{w|u}/2}$, où $q=\operatorname{card}k$, $0\leq w\leq2i$.}

7.4. On peut préciser partiellement ces conjectures, en \unsure{tenant|donnant} compte des représentations de tout le groupe algébrique $\mathbb G_{a\,\mathbb Q_\ell}\times\Gamma_{\mathbb Q_\ell}$, pas seulement de $\Gamma_{\mathbb Q_\ell}$,
